% 1.2 — Point de départ et contribution
\section{Point de départ et contribution}
\label{sec:depart}

\subsection{La version bêta remise au début du stage}
Le stage ne part pas d'une page blanche. WePwn avait publié \demo{} en version
\betaversion{} (\betadate{})~\cite{demo-cli-beta}. Cette bêta posait déjà l'idée qui
reste au c\oe ur du projet, et qu'il faut attribuer à son point d'origine :

\begin{encadre}[L'invariant hérité de la bêta]
Une action qui modifie quelque chose doit être récupérable \emph{et} conforme au contexte
déclaré ; sinon elle est escaladée --- jamais autorisée en silence, jamais présentée à
tort comme « récupérée ».
\end{encadre}

De cet invariant découlait une posture inhabituelle pour un outil de sécurité :
\emph{récupérer plutôt que bloquer}. Quand \demo{} peut prouver qu'il a copié la cible
réelle d'une commande destructrice, il la copie et laisse l'agent continuer ; une
commande \cmd{demo\_cli undo <id>} remet les fichiers en place. Il ne bloque que ce qui
ne peut pas être récupéré (un \cmd{terraform destroy}, un \cmd{git push --force}). La bêta
s'intégrait à Claude Code et à Cursor, savait prévisualiser une requête SQL destructrice
et laissait de chaque décision un reçu chaîné.

Elle annonçait aussi, honnêtement, ses limites. Son journal des modifications le dit en
une phrase : l'analyse des commandes shell repose sur des motifs, « un arbre syntaxique
complet est sur la feuille de route ». Concrètement, le classifieur tenait en
\betaclassify{} lignes et \betaregex{} expressions régulières. Le tableau~\ref{tab:depart}
compare cette bêta à l'état du projet au commit de référence \commitref{}.

\begin{table}[htbp]
\centering
\caption{De la version bêta à l'état actuel (commit \commitref{}, \datecommitref{})}
\label{tab:depart}
\small
\begin{tabularx}{\linewidth}{@{}L{4.1cm}YY@{}}
\toprule
& \textbf{Bêta \betaversion{}} & \textbf{Aujourd'hui} \\
\midrule
Code source (Python) & $\approx$ \betalignessrc{} lignes & $\approx$ \lignessrc{} lignes \\
Tests & $\approx$ \betalignestests{} lignes, 11 fichiers & $\approx$ \lignestests{} lignes,
  \fichierstests{} fichiers \\
Lecture d'une commande & \betaregex{} expressions régulières & analyseur POSIX, simulation sur
  un système de fichiers virtuel, table de \nbregles{} règles (\nboutils{} outils) \\
Décisions & \betadecisions{} & 6 (+ SUPERVISED), canal « demander / refuser » \\
Commandes illisibles & motifs, au mieux & « non compris » : filet, point de contrôle ou
  escalade --- jamais un ALLOW \\
Pendant l'exécution & rien & garde d'appels système (Linux), montage WinFsp (Windows) \\
Surfaces distantes & quelques motifs & $\approx$ \nbsurfaces{} surfaces, chacune avec sa
  règle « ask » ou « deny » \\
Agents & Claude Code, Cursor & Claude Code (Bash, PowerShell, fichiers), Codex, Cursor,
  garde du shell \\
Fichiers de \demo{} & non protégés & plan de contrôle refusé à l'agent \\
Validation & tests unitaires & tests + exemples de règles + portail de \nbcorpus{}
  commandes, Linux et Windows \\
\bottomrule
\end{tabularx}
\end{table}

\subsection{Ce que le stage a apporté}
La mission confiée était de reprendre cette bêta, de la poursuivre et d'en reprendre
l'architecture, en menant de front le développement et le travail offensif. Le dépôt de
travail conserve le détail des \nbcommits{} commits de la dernière phase ; chacun
correspond à une partie nommée, et son message en fait le bilan.

La contribution se résume en quatre lignes :
\begin{enumerate}
  \item \textbf{Un moteur au lieu de motifs} : une analyse syntaxique unique des lignes POSIX,
    une simulation de leurs effets et une table de règles écrite comme des données
    (chapitres~\ref{ch:conception} et~\ref{ch:realisation}).
  \item \textbf{Une seule politique, un seul décideur} : ce qui était dispersé entre
    plusieurs modules est réuni dans \fichier{policy.py} et \fichier{decide.py}.
  \item \textbf{Les deux systèmes} : chaque partie est conçue et vérifiée sous Linux et sous
    Windows ; le filet d'exécution existe sur les deux.
  \item \textbf{Le travail offensif} : un audit pinçant 18 défauts, puis une suite de
    vulnérabilités (V1--V20) trouvées et corrigées, chacune avec un test qui échoue sur le
    code d'avant (chapitre~\ref{ch:securite}).
\end{enumerate}

Une part essentielle du rôle a été celle d'\emph{arbitre}. Chaque cas de politique
délicat --- faut-il autoriser un \cmd{git reset} que le \textit{reflog} sait annuler ?
demander ou refuser la suppression d'un volume Docker ? --- a fait l'objet d'une décision
écrite, numérotée et datée avant d'être codée. Les messages de commit et les tests en
gardent la trace : « ruling (c) », « rulings M1--M5 », « Nizar's table ».
